% Save as: module3-problem-set.tex
% ============================================================================
% Problem Set 3: Producer Theory, Costs & Technology
% Microeconomic Theory for Experimentalists & Applied Microeconomists
% Built from templates/problem_set_template.tex. Compile with pdflatex.
% ============================================================================
\documentclass[11pt]{article}
\usepackage{amsmath, amssymb, amsthm}
\usepackage[margin=1in]{geometry}
\usepackage{enumitem}

\newtheorem{question}{Question}

\title{Problem Set 3: Producer Theory, Costs \& Technology\\
\large Microeconomic Theory for Experimentalists \& Applied Microeconomists}
\author{Due: \textbf{[date --- before Lecture 1 of Module 4]}}
\date{}

\begin{document}
\maketitle

\noindent\emph{Ground rules: collaboration on ideas is encouraged; write-ups
are individual. You may use AI tools for parts flagged with
$\langle$AI$\rangle$ and must attach the transcript. Show all calculations.}

\medskip

\noindent\textbf{Weights:} Part A --- 30 points. Part B --- 45 points.
Part C --- 25 points.

\bigskip

% ---------------------------------------------------------------------------
\section*{Part A: Theory --- Guided Calculations (30 points)}

\begin{question}[Cost minimization and cost curves --- 15 points]
A firm's technology is $q = \sqrt{KL}$. The wage is $w = 4$ and the
capital rental rate is $r = 1$. Inputs are continuously divisible.

\begin{enumerate}[label=(\alph*)]
  \item \textbf{(4 pts)} At the cost-minimizing input mix, the marginal
        rate of technical substitution $MP_L / MP_K = K/L$ equals the
        input price ratio $w/r$. Use this and the production target
        $q = 10$ to show the firm hires $L = 5$ and $K = 20$, at cost 40.
  \item \textbf{(4 pts)} Repeat the logic for a general output target $q$
        to show the long-run cost function is $C(q) = 4q$. What are
        long-run average and marginal cost? In one sentence, connect the
        shape of $C(q)$ to this technology's returns to scale.
  \item \textbf{(5 pts)} Short run: capital is fixed at $\bar K = 20$.
        Show that short-run cost is $SRC(q) = q^2/5 + 20$, that short-run
        average cost $SRAC = q/5 + 20/q$ is minimized at exactly $q = 10$
        with $SRAC = 4$, and that short-run marginal cost equals $SRAC$
        at that point. \emph{Hint: to minimize $SRAC$, set its derivative
        to zero --- or use the fact that MC cuts AC at AC's minimum.}
        State in one sentence what this illustrates about the relation
        between short-run and long-run cost curves.
  \item \textbf{(2 pts)} An artefactual experiment hands manager-subjects
        this exact technology and these prices and asks them to produce
        $q = 10$ (all subjects reach the output target). Subjects' average
        spending is 46, falling to 41 by the tenth round. In two sentences: what economic quantity does the
        gap from 40 measure, and why does its decline with feedback
        matter for interpreting field productivity dispersion?
\end{enumerate}
\end{question}

\begin{question}[Technology adoption as an NPV decision --- 15 points]
Fertilizer costs \$10 per acre, paid at planting. It raises harvest
revenue by \$15 per acre, received one season after planting. A farmer is
deciding whether to adopt on one acre.

\begin{enumerate}[label=(\alph*)]
  \item \textbf{(4 pts)} Compute the NPV of adoption for a farmer whose
        seasonal discount rate is 10\%. Then compute the break-even
        seasonal discount rate --- the rate at which NPV is exactly zero
        --- and say in one sentence why this number makes low adoption a
        puzzle.
  \item \textbf{(4 pts)} Credit: the farmer has no cash at planting and
        must borrow the \$10 at seasonal interest rate $b$, repaying at
        harvest. For $b = 30\%$: does adoption still pay? Find the
        break-even borrowing rate. What does the comparison of the two
        break-even rates in (a) and (b) tell you about how much of the
        adoption puzzle credit constraints alone can explain?
  \item \textbf{(4 pts)} Risk: harvest gains are actually $+\$25$ in a
        good year and $+\$5$ in a bad year, each with probability
        $1/2$ (same mean, \$15). A risk-neutral farmer's decision is
        unchanged --- say why in one sentence. Now consider a
        ``safety-first'' farmer who refuses any option with a possible
        net loss: what does she do, and what intervention does this
        friction respond to that credit and information frictions do not?
  \item \textbf{(3 pts)} A survey finds non-adopting farmers who can
        state the fertilizer return correctly and who hold \emph{liquid}
        savings exceeding \$10 at planting. Which frictions does this observation
        rule out, which remain, and what is the cheapest experiment that
        would test the leading remaining candidate? \emph{(Hint: think
        about WHEN the purchase decision is made.)}
\end{enumerate}
\end{question}

% ---------------------------------------------------------------------------
\section*{Part B: Experimental Design (45 points)}

\begin{question}[Separating the adoption frictions]
Design a field experiment in a smallholder-agriculture setting whose goal
is to \emph{separate} at least three of the four adoption frictions ---
credit constraints, information/learning, risk, and present
bias/procrastination --- for a positive-NPV input (fertilizer or
equivalent). A price-subsidy arm alone is not acceptable: as discussed in
class, a price cut raises adoption under every friction. Specify:
\begin{enumerate}[label=(\alph*)]
  \item \textbf{(6 pts)} Setting, subject pool, and recruitment, and why
        the input's return is credibly positive for this population (your
        design should say what evidence anchors the NPV premise).
  \item \textbf{(12 pts)} Treatment arms. For each arm, state which
        friction it targets and --- the key requirement --- which
        \emph{pattern of effects across arms} identifies each friction.
        Include at least one timing manipulation in the spirit of Duflo,
        Kremer \& Robinson (2011). Say explicitly what the DKR-style
        comparison (small early discount vs.\ large later subsidy)
        identifies that neither arm alone does.
  \item \textbf{(9 pts)} Outcomes and measurement: what you measure, from
        what source (administrative purchase records vs.\ self-reports),
        at what times, and one measurement threat with your answer to it.
  \item \textbf{(9 pts)} Hypotheses and rough sample-size reasoning:
        state the adoption rates you expect by arm under each friction
        story (anchor the control-arm rate and the plausible effect sizes
        on the DKR setting), and the approximate $n$ per arm.
        Simulation-based reasoning is welcome; formulas are not required.
  \item \textbf{(9 pts)} The biggest confound in YOUR design and your
        answer to it. \emph{(Candidates: seasonal liquidity correlates
        with nearly everything; treatment arms that change the number of
        contacts with enumerators; spillovers between neighbors assigned
        to different arms.)}
\end{enumerate}
\end{question}

% ---------------------------------------------------------------------------
\section*{Part C: Silicon Pilot Interpretation $\langle$AI$\rangle$
(25 points)}

\begin{question}[The managers that always adopt]
You run \texttt{module3-simulation-lab.ipynb}: LLM plant managers decide
on the \$10-at-planting / \$15-at-harvest input under three personas
(patient planner, impatient procrastinator, cautious skeptic) and two
decision timings (cash-in-hand vs.\ cash-tight). Adoption comes back at
\textbf{roughly 96\% in every cell}: the timing manipulation moves it by
one percentage point, and the ``impatient procrastinator'' adopts as readily
as the ``patient planner.'' In the human literature, baseline adoption is
a minority and DKR's timing manipulation moved it substantially.

\begin{enumerate}[label=(\alph*)]
  \item \textbf{(8 pts)} Give two mechanisms that could produce this
        pattern, at least one specific to LLMs. \emph{(Think about what a
        transparently positive NPV calculation does to a model trained to
        be helpful and numerate, and about whether a stated persona can
        outweigh an explicit arithmetic cue.)}
  \item \textbf{(9 pts)} Design the diagnostic simulation: propose at
        least two prompt-level manipulations (e.g.\ removing the explicit
        numbers and describing returns qualitatively; adding a concrete
        competing use for the \$10 such as school fees due this week;
        burying the decision in a longer narrative) and state what result
        pattern under each manipulation would tell you. Run it if API
        budget allows and report; otherwise state the predicted patterns
        precisely.
  \item \textbf{(8 pts)} The research group planned to use this sim to
        pick which treatment arms deserve budget in the human experiment
        (e.g.\ ``timing didn't matter in the sim, so drop the DKR arm'').
        Explain why that inference is unsafe, state what the sim CAN
        still safely inform, and give the design change this implies for
        the human pilot plan.
\end{enumerate}
Attach your AI transcripts and, if you ran code, the results CSV.
\end{question}

\end{document}
